% Part: first-order-logic
% Chapter: sequent-calculus
% Section: proving-things-quant

\documentclass[../../../include/open-logic-section]{subfiles}

\begin{document}

\olfileid{fol}{seq}{prq}

\olsection{له سورونو سره \usetoken{P}{derivation}}

\begin{ex}
د سېکوېنټ $\lexists[x][\lnot !A(x)]
\Sequent \lnot \lforall[x][!A(x)]$ يو
$\Log{LK}$-!!{derivation} ورکړئ۔

کله چې له سورونو سره کار کوو، بايد ډاډمن شو چې د ځانګړي متغير شرط
نۀ ماتوو، او کله کله د دې دپاره اړ يو چې د ځينو استنتاجونو د پلي
کولو ترتيب بدل کړو۔ په عمومي توګه دا ګټوره ده چې لومړے هغه قاعدې
وڅېړو چې د ځانګړي متغير د شرط تابع دي (په بشپړ شوي ثبوت کښې به
لاندې وي)۔ همدارنګه، غوره ده لږ وړاندې وګورو او اټکل وکړو چې
لومړنے سېکوېنټ به څنګه وي۔ زموږ په حالت کښې بايد د
$!A(a) \Sequent !A(a)$ په شان وي۔ يعنې کله چې د سورونو قاعدې
“سرچپه” کوو، بايد د $\lforall$ او $\lexists$ دواړو قاعدو دپاره
هماغه اصطلاح---چې $a$ به ئې وبولو---غوره کړو۔ که د هرې قاعدې
دپاره بېله اصطلاح غوره کړو، د $!A(a) \Sequent !A(b)$ په شان څه
ته رسېږو، چې البته د اشتقاق وړ نۀ دے۔

د معمول په شان په دې ليکلو پيل کوو:
\begin{prooftree}
\AxiomC{}
\UnaryInf$\lexists[x][\lnot !A(x)] \fCenter \lnot \lforall[x][!A(x)]$
\end{prooftree}
يا د \LeftR{\exists} قاعده پلې کولے شو او يا د
\RightR{\lnot} قاعده۔ ځکه چې \LeftR{\exists} د ځانګړي متغير د
شرط تابع ده، ښه ده چې دا ژر تر ژره وڅېړو؛ نو لومړے همدا پلې کوو۔
\begin{prooftree}
\AxiomC{}
\UnaryInf$ \lnot !A(a) \fCenter \lnot \lforall[x][!A(x)]$
\RightLabel{\LeftR{\lexists}}
\UnaryInf$ \lexists[x][\lnot !A(x)] \fCenter \lnot \lforall[x][!A(x)]$
\end{prooftree}
د \LeftR{\lnot} او \RightR{\lnot} قاعدو په سرچپه پلي کولو ترلاسه
کوو:
\begin{prooftree}
\AxiomC{}
\UnaryInf$\lforall[x][!A(x)] \fCenter !A(a)$
\RightLabel{\LeftR{\lnot}}
\UnaryInf$\lnot !A(a), \lforall[x][!A(x)] \fCenter $
\RightLabel{\LeftR{\Exchange}}
\UnaryInf$\lforall[x][!A(x)], \lnot !A(a) \fCenter $
\RightLabel{\RightR{\lnot}}
\UnaryInf$ \lnot !A(a) \fCenter \lnot \lforall[x] !A(x)$
\RightLabel{\LeftR{\lexists}}
\UnaryInf$ \lexists[x] \lnot !A(x) \fCenter \lnot \lforall[x] !A(x)$
\end{prooftree}
په دې ځاے کښې يوازينے امکان د \LeftR{\forall} قاعدې پلي کول دي۔
ځکه چې دا قاعده د ځانګړي متغير د بنديز تابع نۀ ده، نو کومه ستونزه
نشته۔ راياد کړئ چې غواړو يو لومړنے سېکوېنټ، يعنې د
$!A(a) \Sequent !A(a)$ په بڼه، ترلاسه کړو؛ نو د قاعدې د کارولو پر
مهال بايد $a$ د~$!A$ د ارګومېنټ په توګه غوره کړو۔
\begin{prooftree}
\Axiom$!A(a) \fCenter !A(a)$
\RightLabel{\LeftR{\lforall}}
\UnaryInf$\lforall[x][!A(x)] \fCenter !A(a)$
\RightLabel{\LeftR{\lnot}}
\UnaryInf$\lnot !A(a), \lforall[x][!A(x)] \fCenter $
\RightLabel{\LeftR{\Exchange}}
\UnaryInf$\lforall[x][!A(x)], \lnot !A(a) \fCenter $
\RightLabel{\RightR{\lnot}}
\UnaryInf$ \lnot !A(a) \fCenter \lnot \lforall[x][!A(x)]$
\RightLabel{\LeftR{\lexists}}
\UnaryInf$ \lexists[x][ \lnot !A(x)] \fCenter \lnot \lforall[x][!A(x)]$
\end{prooftree}
په تېره بيا له سورونو سره د کار پر مهال مهمه ده چې په دې ځاے کښې
بيا وګورو چې د ځانګړي متغير شرط خو نۀ دے مات شوے۔ ځکه چې زموږ له
کارول شوو قاعدو څخه يوازې \LeftR{\exists} د ځانګړي متغير د شرط
تابع وه، او ځانګړے متغير~$a$ د هغې په لاندې سېکوېنټ (وروستي
سېکوېنټ) کښې نۀ راځي، نو دا سم !!{derivation} دے۔
\end{ex}

\begin{prob}
د لاندې سېکوېنټونو !!{derivation}s ورکړئ:
\begin{enumerate}
\item $\Sequent (\lforall[x][!A(x)] \land \lforall[y][!B(y)]) \lif
\lforall[z][(!A(z) \land !B(z))]$.
\item $\Sequent (\lexists[x][!A(x)] \lor \lexists[y][!B(y)]) \lif
\lexists[z][(!A(z) \lor !B(z))]$.
\item $\lforall[x][(!A(x) \lif !B)] \Sequent \lexists[y][!A(y)] \lif !B$.
\item $\lforall[x][\lnot !A(x)] \Sequent \lnot\lexists[x][!A(x)]$.
\item $\Sequent \lnot\lexists[x][!A(x)] \lif \lforall[x][\lnot !A(x)]$.
\item $\Sequent \lnot\lexists[x][\lforall[y][((!A(x,y) \lif \lnot
!A(y,y)) \land (\lnot !A(y,y) \lif !A(x,y)))]]$.
\end{enumerate}
\end{prob}

\begin{prob}
د لاندې سېکوېنټونو !!{derivation}s ورکړئ:
\begin{enumerate}
\item $\Sequent \lnot\lforall[x][!A(x)] \lif \lexists[x][\lnot!A(x)]$.
\item $(\lforall[x][!A(x)] \lif !B) \Sequent \lexists[y][(!A(y) \lif !B)]$.
\item $\Sequent \lexists[x][(!A(x) \lif \lforall[y][!A(y)])]$.
\end{enumerate}
(دا ټول د~$\RightR{\Contraction}$ قاعدې ته اړتيا لري۔)
\end{prob}

\end{document}
